\section{Sketch (Xóa đi khi đã hoàn thiện bài báo)}
Bài báo này sử dụng lối văn phong trung lập, không cảm xúc thiên vị, chủ quan, thẳng thắn nhìn nhận khiếm khuyết của bài báo.

Những thứ sau cần đưa vào bài báo:

\textbf{abstract:} Đề cập đến việc giới thiệu một mô hình toán học dựa trên việc rời rạc hóa không gian tính toán thành các phần tử 3 chiều, sử dụng để mô hình hóa và giải quyết bài toán điện từ cần xét đến cả 3 chiều như máy điện dọc trục; Đề cập đến việc mô hình sử dụng phương pháp từ thế nút, rất dễ thiết lập, có độ tổng quát tốt khi không phải giả thiết trước đường đi của ống từ thông như các phương pháp MEC truyền thống. Tốc độ tính toán có thể nhanh hơn FEM khi kích thước phương trình và độ thưa đều nhỏ hơn FEM. Đề cập đến việc một bài kiểm tra, có thay đổi tham số được thực hiện, cho thấy phương pháp này có thể không ổn định với các bài toán bão hòa sâu.

\textbf{introduction}: 
Bao gồm việc "dẫn truyện" đến phương pháp MEC: là phương pháp cổ điển, tốc độ tính toán cực nhanh, thích hợp khi cần mô phỏng nhanh, quét tham số, tuy nhiên yêu cầu giả thiết trước các ống từ thông (chủ quan). Trái lại, phương pháp phần tử hữu hạn (FEM) được sử dụng rộng rãi và trở thành tiêu chuẩn trong trong nhiều lĩnh vực do độ chính xác cao và tính tổng quát. Phương pháp mạng từ trở dựa trên lưới (MBGRN) là một phương pháp kết hợp được tốc độ tính toán nhanh của MEC và độ chính xác của FEM, (trích dẫn các bài báo đã làm điều này). Sau đó đưa ra vấn đề: Chưa có bài báo nào làm mạng từ trở sinh lưới 3D ( vốn rất quan trọng nếu muốn miêu tả quy luật phân bố từ thông trong không gian 3 chiều, đặc biệt quan trọng với loại động cơ từ thông dọc trục).

Đề cập đến việc phương pháp từ thế nút đơn giản hơn về mặt lập trình, do sự đơn giản hơn trong việc xác định số lượng nút độc lập

Đề cập đến việc sử dụng một lưới có cấu trúc đều theo phương $\theta$, giúp dễ dàng mô phỏng chuyển động quay bằng phép dịch chỉ số các phần tử ở miền quay, do đó tránh được những nhu cầu về tái tạo lại lưới (re-meshing).

Đề cập đến việc phương pháp của tôi có thể cho tốc độ tính toán rất cao, hoặc dung lượng bộ nhớ thấp hơn nhiều so với FEM, do cấu trúc ma trận thưa hơn khoảng từ 3 đến 5 lần so với FEM, và kích thước ma trận cũng nhỏ hơn 

\clearpage